\documentclass[10pt,a4paper]{amsart}
\usepackage{amsmath,amssymb,amsthm,graphicx}
\usepackage{verbatim}
\newenvironment{reference}{\comment}{\endcomment}
%\newenvironment{reference}{}{}
\newenvironment{slogan}{\comment}{\endcomment}
\newenvironment{history}{\comment}{\endcomment}

% For commutative diagrams we use Xy-pic
\usepackage[all]{xy}

% We use 2cell for 2-commutative diagrams.
\xyoption{2cell}
% Theorem environments.
%
\theoremstyle{plain}
\newtheorem{theorem}[subsection]{Theorem}
\newtheorem{proposition}[subsection]{Proposition}
\newtheorem{lemma}[subsection]{Lemma}

\theoremstyle{definition}
\newtheorem{definition}[subsection]{Definition}
\newtheorem{example}[subsection]{Example}
\newtheorem{exercise}[subsection]{Exercise}
\newtheorem{situation}[subsection]{Situation}

\theoremstyle{remark}
\newtheorem{remark}[subsection]{Remark}
\newtheorem{remarks}[subsection]{Remarks}

\numberwithin{equation}{subsection}

% Macros
%
\def\lim{\mathop{\mathrm{lim}}\nolimits}
\def\colim{\mathop{\mathrm{colim}}\nolimits}
\def\Spec{\mathop{\mathrm{Spec}}}
\def\Hom{\mathop{\mathrm{Hom}}\nolimits}
\def\Ext{\mathop{\mathrm{Ext}}\nolimits}
\def\SheafHom{\mathop{\mathcal{H}\!\mathit{om}}\nolimits}
\def\SheafExt{\mathop{\mathcal{E}\!\mathit{xt}}\nolimits}
\def\Sch{\mathit{Sch}}
\def\Mor{\mathop{\mathrm{Mor}}\nolimits}
\def\Ob{\mathop{\mathrm{Ob}}\nolimits}
\def\Sh{\mathop{\mathit{Sh}}\nolimits}
\def\NL{\mathop{N\!L}\nolimits}
\def\CH{\mathop{\mathrm{CH}}\nolimits}
\def\proetale{{pro\text{-}\acute{e}tale}}
\def\etale{{\acute{e}tale}}
\def\QCoh{\mathit{QCoh}}
\def\Ker{\mathop{\mathrm{Ker}}}
\def\Im{\mathop{\mathrm{Im}}}
\def\Coker{\mathop{\mathrm{Coker}}}
\def\Coim{\mathop{\mathrm{Coim}}}

% Boxtimes
%
\DeclareMathSymbol{\boxtimes}{\mathbin}{AMSa}{"02}

%
% Macros for moduli stacks/spaces
%
\def\QCohstack{\mathcal{QC}\!\mathit{oh}}
\def\Cohstack{\mathcal{C}\!\mathit{oh}}
\def\Spacesstack{\mathcal{S}\!\mathit{paces}}
\def\Quotfunctor{\mathrm{Quot}}
\def\Hilbfunctor{\mathrm{Hilb}}
\def\Curvesstack{\mathcal{C}\!\mathit{urves}}
\def\Polarizedstack{\mathcal{P}\!\mathit{olarized}}
\def\Complexesstack{\mathcal{C}\!\mathit{omplexes}}
% \Pic is the operator that assigns to X its picard group, usage \Pic(X)
% \Picardstack_{X/B} denotes the Picard stack of X over B
% \Picardfunctor_{X/B} denotes the Picard functor of X over B
\def\Pic{\mathop{\mathrm{Pic}}\nolimits}
\def\Picardstack{\mathcal{P}\!\mathit{ic}}
\def\Picardfunctor{\mathrm{Pic}}
\def\Deformationcategory{\mathcal{D}\!\mathit{ef}}
\usepackage{cleveref}
\crefname{originallemma}{original lemma}{original lemmas}
\Crefname{originallemma}{Original lemma}{Original lemmas}
\makeatletter
\newlabel{morphisms-lemma-affine-equivalence-modules}{{01SB}{source}{}{sourcepointer-1}{}}
\newlabel{morphisms-lemma-affine-equivalence-modules@cref}{{[originallemma][0][]01SB}{[1][0][]source}{}{}{}}
\makeatother
\crefname{originalreference}{original reference}{original references}
\Crefname{originalreference}{Original reference}{Original references}
\makeatletter
\newlabel{proposition-projective-line}{{0C6U}{source}{}{sourcepointer-2}{}}
\newlabel{proposition-projective-line@cref}{{[originalreference][0][]0C6U}{[1][0][]source}{}{}{}}
\makeatother
\crefname{originallemma}{original lemma}{original lemmas}
\Crefname{originallemma}{Original lemma}{Original lemmas}
\makeatletter
\newlabel{varieties-lemma-divisible}{{0C6P}{source}{}{sourcepointer-3}{}}
\newlabel{varieties-lemma-divisible@cref}{{[originallemma][0][]0C6P}{[1][0][]source}{}{}{}}
\makeatother
\crefname{originallemma}{original lemma}{original lemmas}
\Crefname{originallemma}{Original lemma}{Original lemmas}
\makeatletter
\newlabel{varieties-lemma-regular-functions-proper-variety}{{04L2}{source}{}{sourcepointer-4}{}}
\newlabel{varieties-lemma-regular-functions-proper-variety@cref}{{[originallemma][0][]04L2}{[1][0][]source}{}{}{}}
\makeatother
\crefname{originalsection}{original section}{original sections}
\Crefname{originalsection}{Original section}{Original sections}
\makeatletter
\newlabel{varieties-section-normalization}{{0BXQ}{source}{}{sourcepointer-5}{}}
\newlabel{varieties-section-normalization@cref}{{[originalsection][0][]0BXQ}{[1][0][]source}{}{}{}}
\makeatother
\crefname{originalsection}{original section}{original sections}
\Crefname{originalsection}{Original section}{Original sections}
\makeatletter
\newlabel{varieties-section-normalization-one-dimensional}{{0C44}{source}{}{sourcepointer-6}{}}
\newlabel{varieties-section-normalization-one-dimensional@cref}{{[originalsection][0][]0C44}{[1][0][]source}{}{}{}}
\makeatother
\title{A singular genus 0 proper curve with H\textasciicircum{}0(X, O\_X) = k is P\textasciicircum{}1 over a nontrivial finite extension glued to Spec k at one rational point}
\author{The Stacks Project Authors}
\date{stacks-project:curves.tex:local-snapshot-20261012; GNU FDL 1.2 or later}
\begin{document}
\maketitle
\noindent\textit{Editorial note.} All delivered mathematical text is the author's own, copied from the pinned Stacks Project chapter file curves.tex; no statement or proof step was rewritten, reordered or removed.\par\medskip
\noindent\textit{Editorial note.} The counted proof is the author's complete argument: the normalization sequence, the conductor ideal with its defining exact sequence, the commutative-subalgebra dimension bound, the descent I = nu\_*I', the estimate dim H\textasciicircum{}0(X', A') <= 2n-2, the divisibility argument, and the final identification of both rational points and the structure sequence are all retained verbatim.\par\medskip
\noindent\textit{Editorial note.} The label lemma-genus-zero-singular names the counted statement, and the proof environment immediately following it in the source is delivered verbatim as its proof.\par\medskip
\noindent\textit{Editorial note.} The counted statement is carried by the author's own lemma environment in the chapter source.\par\medskip
\noindent\textit{Allowed standard prerequisites (editorial).}\begin{itemize}
\item Normalization of a singular curve. Varieties, Sections varieties-section-normalization and varieties-section-normalization-one-dimensional, with Lemma varieties-lemma-regular-functions-proper-variety. Supplies the finite normalization nu : X' -> X, the finite extension k' = H\textasciicircum{}0(X', O\_\{X'\}) and the short exact sequence 0 -> O\_X -> nu\_*O\_\{X'\} -> Q -> 0 supported in finitely many points.
\item Classification of genus 0 curves without H\textasciicircum{}0 correction. Same chapter, Proposition proposition-projective-line. Identifies X' as P\textasciicircum{}1\_\{k'\} once H\textasciicircum{}1(X', O\_\{X'\}) = 0 and the conductor quotient is trivial on X'.
\end{itemize}\medskip
\ifdefined\hypertarget\hypertarget{sourcepointer-1}{}\fi
\noindent\textit{Source reference.} Original lemma 01SB: sources/stacks-chow/tags.txt, source lines 2332--2332. Prerequisite pointer invoked at proof line(s) 2301; it names a standard fact of a sibling chapter that the counted argument uses, and every step of the counted proof built on it is retained verbatim in the delivered body.\par\medskip
\ifdefined\hypertarget\hypertarget{sourcepointer-2}{}\fi
\noindent\textit{Source reference.} Original reference 0C6U: sources/stacks-curves/pinned-original-source.tex, source lines 2138--2138. Prerequisite pointer invoked at proof line(s) 2320; it names a standard fact of the same chapter that the counted argument uses, and every step of the counted proof built on it is retained verbatim in the delivered body.\par\medskip
\ifdefined\hypertarget\hypertarget{sourcepointer-3}{}\fi
\noindent\textit{Source reference.} Original lemma 0C6P: sources/stacks-chow/tags.txt, source lines 15304--15304. Prerequisite pointer invoked at proof line(s) 2316; it names a standard fact of a sibling chapter that the counted argument uses, and every step of the counted proof built on it is retained verbatim in the delivered body.\par\medskip
\ifdefined\hypertarget\hypertarget{sourcepointer-4}{}\fi
\noindent\textit{Source reference.} Original lemma 04L2: sources/stacks-chow/tags.txt, source lines 5993--5993. Prerequisite pointer invoked at proof line(s) 2251; it names a standard fact of a sibling chapter that the counted argument uses, and every step of the counted proof built on it is retained verbatim in the delivered body.\par\medskip
\ifdefined\hypertarget\hypertarget{sourcepointer-5}{}\fi
\noindent\textit{Source reference.} Original section 0BXQ: sources/stacks-chow/tags.txt, source lines 14988--14988. Prerequisite pointer invoked at proof line(s) 2247; it names a standard fact of a sibling chapter that the counted argument uses, and every step of the counted proof built on it is retained verbatim in the delivered body.\par\medskip
\ifdefined\hypertarget\hypertarget{sourcepointer-6}{}\fi
\noindent\textit{Source reference.} Original section 0C44: sources/stacks-chow/tags.txt, source lines 15214--15214. Prerequisite pointer invoked at proof line(s) 2248; it names a standard fact of a sibling chapter that the counted argument uses, and every step of the counted proof built on it is retained verbatim in the delivered body.\par\medskip
\begin{lemma}
\label{lemma-genus-zero-singular}
Let $X$ be a proper curve over a field $k$ with $H^0(X, \mathcal{O}_X) = k$.
Assume $X$ is singular and has genus $0$. Then there exists a diagram
$$
\xymatrix{
x' \ar[d] \ar[r] & X' \ar[d]^\nu \ar[r] & \Spec(k') \ar[d] \\
x \ar[r] & X \ar[r] & \Spec(k)
}
$$
where
\begin{enumerate}
\item $k'/k$ is a nontrivial finite extension,
\item $X' \cong \mathbf{P}^1_{k'}$,
\item $x'$ is a $k'$-rational point of $X'$,
\item $x$ is a $k$-rational point of $X$,
\item $X' \setminus \{x'\} \to X \setminus \{x\}$ is an isomorphism,
\item $0 \to \mathcal{O}_X \to \nu_*\mathcal{O}_{X'} \to k'/k \to 0$
is a short exact sequence
where $k'/k = \kappa(x')/\kappa(x)$ indicates the skyscraper sheaf
on the point $x$.
\end{enumerate}
\end{lemma}

\begin{proof}
Let $\nu : X' \to X$ be the normalization of $X$, see
Varieties, Sections \ref{varieties-section-normalization} and
\ref{varieties-section-normalization-one-dimensional}.
Since $X$ is singular $\nu$ is not an isomorphism.
Then $k' = H^0(X', \mathcal{O}_{X'})$ is a finite extension of $k$
(Varieties, Lemma \ref{varieties-lemma-regular-functions-proper-variety}).
The short exact sequence
$$
0 \to \mathcal{O}_X \to \nu_*\mathcal{O}_{X'} \to \mathcal{Q} \to 0
$$
and the fact that $\mathcal{Q}$ is supported in finitely many
closed points give us that
\begin{enumerate}
\item $H^1(X', \mathcal{O}_{X'}) = 0$, i.e., $X'$ has genus $0$
as a curve over $k'$,
\item there is a short exact sequence
$0 \to k \to k' \to H^0(X, \mathcal{Q}) \to 0$.
\end{enumerate}
In particular $k'/k$ is a nontrivial extension.

\medskip\noindent
Next, we consider what is often called the {\it conductor ideal}
$$
\mathcal{I} = \SheafHom_{\mathcal{O}_X}(\nu_*\mathcal{O}_{X'}, \mathcal{O}_X)
$$
This is a quasi-coherent $\mathcal{O}_X$-module. We view $\mathcal{I}$
as an ideal in $\mathcal{O}_X$ via the map $\varphi \mapsto \varphi(1)$.
Thus $\mathcal{I}(U)$ is the set of $f \in \mathcal{O}_X(U)$ such that
$f \left(\nu_*\mathcal{O}_{X'}(U)\right) \subset \mathcal{O}_X(U)$. In
other words, the condition is that $f$ annihilates $\mathcal{Q}$.
In other words, there is a defining exact sequence
$$
0 \to \mathcal{I} \to \mathcal{O}_X \to
\SheafHom_{\mathcal{O}_X}(\mathcal{Q}, \mathcal{Q})
$$
Let $U \subset X$ be an affine open containing the support of $\mathcal{Q}$.
Then $V = \mathcal{Q}(U) = H^0(X, \mathcal{Q})$ is a $k$-vector space
of dimension $n - 1$. The image of
$\mathcal{O}_X(U) \to \Hom_k(V, V)$ is a commutative subalgebra,
hence has dimension $\leq n - 1$ over $k$ (this is a property of
commutative subalgebras of matrix algebras; details omitted).
We conclude that we have a short exact sequence
$$
0 \to \mathcal{I} \to \mathcal{O}_X \to \mathcal{A} \to 0
$$
where $\text{Supp}(\mathcal{A}) = \text{Supp}(\mathcal{Q})$
and $\dim_k H^0(X, \mathcal{A}) \leq n - 1$.
On the other hand, the description
$\mathcal{I} = \SheafHom_{\mathcal{O}_X}(\nu_*\mathcal{O}_{X'}, \mathcal{O}_X)$
provides $\mathcal{I}$ with a $\nu_*\mathcal{O}_{X'}$-module structure
such that the inclusion map $\mathcal{I} \to \nu_*\mathcal{O}_{X'}$
is a $\nu_*\mathcal{O}_{X'}$-module map.
We conclude that $\mathcal{I} = \nu_*\mathcal{I}'$
for some quasi-coherent sheaf of ideals
$\mathcal{I}' \subset \mathcal{O}_{X'}$, see
Morphisms, Lemma \ref{morphisms-lemma-affine-equivalence-modules}.
Define $\mathcal{A}'$ as the cokernel:
$$
0 \to \mathcal{I}' \to \mathcal{O}_{X'} \to \mathcal{A}' \to 0
$$
Combining the exact sequences so far we obtain a short exact sequence
$0 \to \mathcal{A} \to \nu_*\mathcal{A}' \to \mathcal{Q} \to 0$.
Using the estimate above,
combined with $\dim_k H^0(X, \mathcal{Q}) = n - 1$, gives
$$
\dim_k H^0(X', \mathcal{A}') =
\dim_k H^0(X, \mathcal{A}) + \dim_k H^0(X, \mathcal{Q}) \leq 2 n - 2
$$
However, since $X'$ is a curve over $k'$ we see that
the left hand side is divisible by $n$
(Varieties, Lemma \ref{varieties-lemma-divisible}).
As $\mathcal{A}$ and $\mathcal{A}'$ cannot be zero, we conclude that
$\dim_k H^0(X', \mathcal{A}') = n$ which means that $\mathcal{I}'$
is the ideal sheaf of a $k'$-rational point $x'$.
By Proposition \ref{proposition-projective-line}
we find $X' \cong \mathbf{P}^1_{k'}$.
Going back to the equalities above, we conclude that
$\dim_k H^0(X, \mathcal{A}) = 1$. This
means that $\mathcal{I}$ is the ideal sheaf of a
$k$-rational point $x$. Then $\mathcal{A} = \kappa(x) = k$
and $\mathcal{A}' = \kappa(x') = k'$ as skyscraper sheaves.
Comparing the exact sequences given above,
this immediately implies the result on structure sheaves
as stated in the lemma.
\end{proof}

\end{document}
